% =====================================================================
%  Chapter 8 - Legal, Social, Ethical and Professional Issues
%  Rewritten 2026-08-05. Adds the two things the rubric names explicitly
%  and the draft omitted: professional standards / codes of conduct, and
%  intellectual property + software trustworthiness. Existing labels kept.
% =====================================================================
\chapter{Legal, Social, Ethical and Professional Considerations}
\label{ch:lsep}

\section{Professional Standards and Codes of Conduct}
\label{sec:lsep-professional}

Two codes of conduct govern work of this kind and both were treated as design
constraints rather than as a compliance exercise appended afterwards. The BCS
Code of Conduct requires due regard for public health and safety, working
within the limits of one's competence, and not misrepresenting the capability
of a product or service, and the ACM Code of Ethics adds an obligation to be
honest about a system's limitations and to give comprehensive evaluations of
its risks.

Three decisions in this project follow directly from those obligations. The
first is that the safety envelope of Section~\ref{sec:eval-envelope} is
reported as empty at 3B and as clearing only the low-risk bar at 7B and 14B,
which is the least flattering way to present the result and the only one
compatible with not misrepresenting capability. The second is the treatment of
the off-by-one defect in the entropy gate's logit slice
(Section~\ref{sec:disc-limitations}), which was found late, could not be
corrected without re-running calibration on all three models, and is therefore
documented with its measurement and its consequences rather than quietly left
in place. The third is that a refuted hypothesis is reported as refuted in the
gate ablation of Section~\ref{sec:eval-ablation}, where the author's stated
prior expectation was that up-weighting the probe would help and the data
showed the opposite.

Competence limits matter here too. The clinical accuracy bars used throughout
this dissertation are stipulated by the author and are not derived from
clinical evidence or endorsed by any clinician, and the risk tagger is a
keyword rule set written by a computer scientist rather than a triage
instrument validated by a clinical team. Both are stated as such wherever they
are used, because presenting a stipulated threshold as a clinical standard
would be exactly the misrepresentation the codes prohibit.

Finally, no human participants and no patient data were involved at any stage,
so the project required no ethics approval under the College's research
governance framework. All data used is publicly available research material,
and the position is stated here rather than assumed.

\section{Patient Safety and Clinical Liability}
\label{sec:lsep-safety}

The asymmetry of harm in clinical question answering, where an unnecessary
retrieval costs a few seconds of latency and a confidently wrong
drug-interaction answer could cost a patient, is encoded in the artefact
rather than only asserted in the write-up. The gate ensemble is deliberately
biased toward retrieving when its signals disagree
(Section~\ref{sec:method-why-three}), and the risk tagger's precedence rules
default toward the more cautious label.

The evaluation then returned a result this chapter has to state as plainly as
the technical chapters do. No policy clears the lowest stipulated accuracy bar
at 3B, and the larger models clear the low-risk bar only. This system must
therefore not be used for clinical decision support, including in
education-adjacent settings, without qualified human oversight. Quantifying
the distance to acceptability is the safety-relevant contribution the project
can honestly make, and concealing it would have been the ethical failure.

Liability follows the same logic. If a system of this kind were deployed and
gave a harmful answer, responsibility would distribute contentiously across
the model provider, the deployer and the clinician who acted on the output,
and regulatory regimes increasingly refuse to let software disclaim its way
out of that. Under the UK MHRA's guidance on software as a medical device, and
under the EU AI Act's classification of clinical decision support as
high-risk, a deployed version of this system would attract conformity
obligations that this research prototype has never been assessed against. It
is a benchmark instrument rather than a device.

One empirical finding sharpens the automation-bias concern beyond the generic
version of it. Section~\ref{sec:eval-distraction} showed that when retrieval
misleads the model, its wrong answers become longer and more fluent, running
roughly $15\times$ the length of the terse correct closed-book answers.
Fluency and apparent thoroughness are precisely the surface features that
induce human over-trust, so the failure mode this system exhibits is the most
dangerous kind for a clinical reader, and that is an independent argument for
surfacing machine uncertainty (Sections~\ref{sec:impl-demo}
and~\ref{sec:disc-deployment}) rather than only polished prose.

\section{Privacy, Data Governance and Sovereignty}
\label{sec:lsep-privacy}

The offline, CPU-only architecture is a privacy architecture as much as a cost
profile. No query, no patient context and no retrieved passage ever leaves the
device, so there is no API call to log, no third-party processor to contract
with and no cross-border transfer to assess. That is what makes the design
viable in precisely the environments where GDPR, NHS data-governance
frameworks and analogous regimes make cloud LLM use difficult or
impermissible, and it is why offline operation appears in the problem
statement rather than among the limitations.

What the project did not have to handle belongs here as well. The corpus is
built from public sources and the benchmarks are public research datasets, so
no patient data was processed at any point, and the credentialed EHR datasets
contemplated in the original project plan were never used. The privacy
property of the architecture is real, but it was demonstrated on public data,
and the accurate claim is an architectural one, which is that a system with no
network egress cannot leak what it processes.

\section{Intellectual Property, Licensing and Software Trustworthiness}
\label{sec:lsep-ip}

Four distinct licences meet in this artefact and they do not grant the same
things. The code written for this project is released under the MIT licence,
and Qwen2.5-7B and 14B are distributed under Apache 2.0, which places no
restriction on the use made of them here. Llama-3.2-3B-Instruct is distributed
under the Llama 3.2 Community Licence, which is not an open-source licence in
the OSI sense, since it carries an acceptable-use policy and attribution
requirements that continue to bind derived work, and academic evaluation of
this kind falls within it. The MedRAG corpora and the MIRAGE and PubMedQA
benchmarks are third-party research datasets used under their own terms, and
none is redistributed in the project archive, which ships only the scripts
that fetch them. The derived FAISS and BM25 indexes are excluded for the same
reason.

Software trustworthiness is a professional obligation as well as an
engineering one, because a system whose numbers cannot be checked cannot
responsibly inform a safety-relevant conclusion. The machinery of
Section~\ref{sec:impl-verification} is the concrete response, since every
reported number is generated from raw run logs, a regression test fails the
build if any value drifts from what the logs produce, and the logs are shipped
so a reader can regenerate every table and figure without downloading a model.
In a domain where an unverifiable claim can propagate into clinical use, that
auditability is an ethical property rather than a technical convenience.

\section{Bias, Equity and Access}
\label{sec:lsep-bias}

Neither of the gate's two paths is neutral. When the system retrieves it
inherits the perspective of its corpus, which is Western-published textbooks
and PubMed abstracts whose epidemiological assumptions, drug availability and
demographic baselines are present in every passage. When the gate skips it
inherits the training distribution of the base model instead. A gate therefore
does not remove bias, it selects between two bias profiles on every query, and
no analysis in this study measures the demographic consequences of that
selection, which is an acknowledged gap and a natural extension of the
risk-stratified evaluation already in place.

The equity argument that motivates the project also carries its own warning.
An offline CPU-only system is exactly the configuration that can run where
GPUs, budgets and connectivity do not exist, in rural practice, LMIC clinics
and training hospitals, the settings where only 14\% of rural residents in
low-income countries are online at all \citep{itu2025factsfigures}. Those
settings have the most to gain from deployable clinical question answering and
the least infrastructure with which to verify it. Since
Section~\ref{sec:eval-envelope} shows the present system does not meet even
lenient safety bars, shipping it to under-resourced settings where clinical
backstop capacity is thinnest would be worse than shipping nothing. Equity of
access has to mean access to a safe system, and this work measures how far
that still is.

Accessibility of the artefact itself is a smaller but real gap. The
demonstration interface of Section~\ref{sec:impl-demo} encodes gate
uncertainty as colour intensity alone, which is not usable by a colour-blind
reader without the numeric signal alongside it, and that is recorded as a
defect to fix rather than defended.

\section{Environmental Cost}
\label{sec:lsep-energy}

Per-query energy was measured throughout with \texttt{codecarbon} and is a
software estimate rather than a metered hardware measurement, and is reported
as such. The honest accounting cuts against the intuitive environmental claim,
because the gate halves retrieval calls but adds three draft generations per
query, so on hardware where generation dominates energy consumption gating is
not currently an energy win (Section~\ref{sec:eval-cost}). The case for
gating on energy grounds is real but conditional, holding where retrieval
dominates the cost profile, as with large corpora or remote and disk-bound
indexes, rather than on this laptop. No unearned environmental benefit is
claimed, and the merged-draft optimisation of Section~\ref{sec:conc-future} is
the identified route to reducing the gate's own footprint. More broadly, the
choice of a small quantised model on commodity hardware sits at the frugal end
of the field's energy spectrum \citep{strubell2019}, which is itself an
environmentally relevant design position.